\section{Sharpness of the uniform primitive modulus}
\label{sec:r80-sharp-stability}

The exponent in Theorem~\ref{thm:r79-frontier-stability} cannot be
increased on its stated class.  This failure of a Lipschitz estimate
already occurs for two scalar models, two commands and two live labels.
The values remain finite, the call costs are identically one, and the
same one-row fallback works throughout the perturbation.  Thus the
phenomenon is not caused by a loss of feasibility or by an increasing
physical dimension.  It comes from the cost of repairing a label which
was dark and unreachable in one of the models.

Put $\beta=1/4$ and $q=1-\beta=3/4$.  Let
$\mathcal U_\varepsilon=\{0_\varepsilon,1\}$, where
$0\le\varepsilon<1$, be scalar instruments with unit call cost,
one failure report, and the following success probabilities:
\begin{equation}\label{eq:r80-sharp-models}
\begin{array}{c|cc}
 &a&b\\ \hline
 0_\varepsilon&1-\varepsilon&0\\
 1&0&\beta.
\end{array}
\end{equation}
The model is fixed throughout the run and is unobserved.  Success
requires immediate stopping; the sampled command is known to the
controller.  For a scalar instrument, the displayed probabilities
specify all the nonterminal maps.  These are valid instruments for
every stated parameter, with positive deterministic paid costs and
no additional observation.

The controller convention is that of
\eqref{eq:r79-controller-parameters}: a prescribed pure initial live
label, stochastic command and update rows, and no free clock or
persistent private seed.  We write $V_2(\varepsilon)$ for its
two-live-label robust optimum.  Define
\begin{equation}\label{eq:r80-sharp-constants}
 r_0=\frac{\sqrt{65}-1}{32},\qquad
 B_0=\frac1{r_0}=\frac{1+\sqrt{65}}2,\qquad
 D=\beta+qr_0,\qquad
 c_0=\frac{2\sqrt q\,r_0^2}{D^2+\beta r_0^2}.
\end{equation}
In particular $c_0>0$; its approximate value is $0.4564191442$.

\begin{theorem}[An exact square-root singularity at fixed memory]
\label{thm:r80-sharp-modulus}
For the primitive family \eqref{eq:r80-sharp-models},
\begin{equation}\label{eq:r80-sharp-asymptotic}
 V_2(0)=B_0,\qquad
 \lim_{\varepsilon\downarrow0}
       \frac{V_2(\varepsilon)-B_0}{\sqrt\varepsilon}=c_0.
\end{equation}
The minimum defining each $V_2(\varepsilon)$ is attained.  The
instrument discrepancy between the paired families
$\mathcal U_0$ and $\mathcal U_\varepsilon$ is exactly
$\varepsilon$, and their cost discrepancy is zero.
For $0\le\varepsilon\le1/2$ they have unit cost floor and ceiling
and a common one-row cost bound $6$.

Consequently no local uniform bound with exponent strictly greater
than $1/2$ can replace the instrument term in
Theorem~\ref{thm:r79-frontier-stability}, even at the single fixed
budget $M=2$.  In contrast, the one-label frontier of the same
family is
\begin{equation}\label{eq:r80-one-label-smooth}
 V_1(\varepsilon)=\frac1{1-\varepsilon}+\frac1\beta.
\end{equation}
The exponent obstruction therefore concerns an optimized memory
frontier, rather than a scalar waiting-time singularity.
\end{theorem}

We prove the theorem below.  Two elementary facts about scalar
two-label chains will be useful.  If the success probabilities at
the labels are $h_1,h_2$ and the probabilities of a nonterminal
transition to the other label are $t_1,t_2$, respectively, the
finite cost from label $1$ is
\begin{equation}\label{eq:r80-two-state-cost}
 C_1=\frac{h_2+t_2+t_1}{h_1h_2+h_1t_2+h_2t_1}.
\end{equation}
This formula applies when the denominator is positive.  When the
second label is inaccessible, its possibly infinite cost must instead
be discarded: the initial cost is $1/h_1$.  These conclusions follow
directly from
$C_1=1+(1-h_1-t_1)C_1+t_1C_2$ and its second-row equation.
Whenever both costs are finite, their ordering is the reverse of
the ordering of $h_1,h_2$.  Indeed
\[
 (h_2-h_1)C_1=(h_2+t_1+t_2)(C_1-C_2).
\]
Thus routing a given failure to the label with larger success
probability cannot increase the cost when that routing does not
affect the other model.  This last statement can also be checked by
substituting the old positive cost vector into the new Bellman
inequality; finite-prefix iteration proves the comparison without
an unproved properness assumption.

\subsection{The unperturbed optimum and its rigidity}

Specify every two-label controller by six numbers in $[0,1]$:
\begin{equation}\label{eq:r80-six-parameters}
\begin{array}{c|l}
r,s&\text{probabilities of command $a$ at labels $1,2$},\\
\alpha,\gamma&\text{probabilities of writing label $2$ after,
respectively, an $a$- or $b$-failure at label $1$},\\
\xi,\delta&\text{probabilities of writing label $1$ after,
respectively, an $a$- or $b$-failure at label $2$}.
\end{array}
\end{equation}
Label $1$ is the initial label.  There is no restriction on any of
these probabilities.  The description includes one-label controllers
by allowing the second label to remain unused.

\begin{lemma}[The unique effective optimal program at zero perturbation]
\label{lem:r80-boundary-rigidity}
At $\varepsilon=0$, every robust optimal two-label controller has
\begin{equation}\label{eq:r80-rigid-controller}
 r=r_0,\qquad s=0,\qquad
 \alpha=1,\qquad\gamma=\delta=0.
\end{equation}
The value of $\xi$ is immaterial because command $a$ is never used
at label $2$.  The optimum is $B_0$.
Every controller with $r\le s$ has robust cost at least
\begin{equation}\label{eq:r80-other-orientation}
 B_{\rm opp}=2(\sqrt{19}-2)>B_0.
\end{equation}
The lower bound in \eqref{eq:r80-other-orientation} is attained
within that orientation.
\end{lemma}
\begin{proof}
In model $0_0$, an $a$ call always succeeds.  Changing updates after
an $a$-failure therefore leaves this model's cost unchanged.

First suppose $r\ge s$.  In model $1$ the second label has at least
as large a success probability as the first.  By the preceding
two-label comparison, routing every $a$-failure to label $2$ can
only improve model $1$.  Hence it gives a lower bound for an
arbitrary controller.  This routing sets $\alpha=1,\xi=0$.
The endpoint $r=0$ makes model $0$ improper.  At $r=1$ the first
call in model $1$ is certainly a failure and at least $1/\beta$
further calls are required in expectation, giving cost at least
$5$.  Both endpoints are therefore harmless for a lower bound
strictly below $5$; in the following division take $0<r<1$.
Let $C_0,C_1$ denote the two initial costs after this improvement.
For their explicit expressions put
\begin{align*}
 N_0&=s+(1-s)\delta+(1-r)\gamma,\\
 D_0&=rs+r(1-s)\delta+s(1-r)\gamma,\\
 N_1&=\beta(1-s)+q(1-s)\delta+r+q(1-r)\gamma,\\
 D_1&=\beta(1-s)
            [\beta+qr+q(1-r)(\delta+\gamma)].
\end{align*}
When the relevant denominators are positive,
$C_0=N_0/D_0$ and $C_1=N_1/D_1$.  Define
\begin{equation}\label{eq:r80-comparison-function}
 f(t)=\frac{\beta+t}{\beta(\beta+qt)},\qquad 0\le t\le1.
\end{equation}
It is strictly increasing.  The inequality
\begin{equation}\label{eq:r80-palette-comparison}
 C_1\ge f(1/C_0)
\end{equation}
has the following explicit positive-polynomial certificate:
\begin{align}
 &N_1\beta(\beta N_0+qD_0)
               -D_1(\beta N_0+D_0) \nonumber\\
 ={}&\beta q s^2(1-r)^2\gamma^2
       +2\beta qsr(1-s)(1-r)\gamma\delta \nonumber\\
 &+\beta(1-r)
      [\beta^2(r-s)+s^2(\beta+2qr)]\gamma \nonumber\\
 &+\beta q r^2(1-s)^2\delta^2
       +\beta sr(1-s)(\beta+2qr)\delta
       +\beta s^2r(\beta+qr).
       \label{eq:r80-positive-polynomial}
\end{align}
Every summand on the right is nonnegative for $0\le s\le r\le1$.
This identity follows by expansion and proves the claimed inequality
after division by its positive denominator.  The formula does not
use the value of $\beta$ beyond $0<\beta<1$.

The inaccessible case omitted by this division is
$s=\delta=\gamma=0$.  A finite initial cost then requires $r>0$.
Direct conditioning gives $C_0=1/r$ and, after the indicated
$a$-failure improvement, $C_1=f(r)$.  If other endpoint choices
make one of the costs infinite, they cannot improve a finite
optimum.  The remaining finite endpoints are covered either by
the displayed formulas or by their one-label reduction.

If both costs are at most $b$, \eqref{eq:r80-palette-comparison}
implies $b\ge f(1/b)$.  The positive solution of
$b=f(1/b)$ is $B_0$, since its reciprocal solves
$r^2+\beta^2r-\beta^2=0$.  This proves the lower bound for
$r\ge s$.  Equality is achieved by
\eqref{eq:r80-rigid-controller}.  At a finite optimum one has
$0<r<1$; the endpoint $r=0$ cannot succeed in model $0$, whereas
$r=1$ gives model $1$ cost at least $1+1/\beta=5>B_0$.
In this open range the right-hand side of
\eqref{eq:r80-positive-polynomial} is strictly positive if any
of $s,\delta,\gamma$ is positive.  Hence equality requires all
three to vanish.  The two remaining inequalities then force
$r=r_0$.  Without the $a$-failure improvement, the model-$1$ cost
at these rows is
\[
 \frac{\beta+r\alpha}
      {\beta[\beta(1-r)+r\alpha]}.
\]
It is strictly decreasing in $\alpha$ for $0<r<1$.
Its value equals $B_0$ only at $\alpha=1$.  This proves the
stated rigidity in this orientation.

It remains to exclude $r\le s$.  In this orientation, routing
every $a$-failure to label $1$ improves model $1$ and leaves
model $0$ unchanged.  Make that improvement.  The same $N_0,D_0$
describe model $0$, whereas model $1$ has
\begin{align*}
 \widetilde N_1&=\beta+qs+q(1-s)\delta+q(1-r)\gamma,\\
 \widetilde D_1&=\beta(1-r)
                [\beta+qs+q(1-s)(\delta+\gamma)].
\end{align*}
Suppose a proposed bound $b<B_{\rm opp}<5$ holds for both
models.  Model $1$ has success probability at most
$\beta(1-r)$ at every call, so
$b\ge1/[\beta(1-r)]$.  In particular $b\ge4$ and
$br<1$, the latter because $b<1+1/\beta=5$.
The condition $N_0/D_0\le b$ now implies $bs>1$ and
\begin{equation}\label{eq:r80-gamma-bound}
 \gamma\ge
 \frac{[s+(1-s)\delta](1-br)}{(1-r)(bs-1)}.
\end{equation}
The model-$1$ cost increases with $\gamma$: it moves a failure
from the better first label to the worse second label.  Substitute
the right-hand side of \eqref{eq:r80-gamma-bound} in
$\widetilde N_1/\widetilde D_1$, and call the resulting rational
function $F_b(r,s,\delta)$.  Its denominator is positive on the
range in question, even if the substituted $\gamma$ exceeds one.
The latter substitution only relaxes the original constraint.

For precision, differentiation gives the following signs:
\begin{align*}
 \partial_\delta F_b(r,s,\delta)&\ge0,\\
 \partial_r F_b(r,s,0)&>0,\\
 \operatorname{sign}\partial_s F_b(0,s,0)
        &=\operatorname{sign}(bs-2).
\end{align*}
These signs can be verified without estimating a root.  Up to
positive squared denominators, their respective numerators are
\[
 (\beta-1)(r-s)(s-1)(br-1)(bs-1),\quad
 -(bs-1)^2(\beta s-\beta-s),\quad
 (1-\beta)s(bs-2).
\]
Consequently
\[
 C_1\ge F_b(r,s,\delta)
       \ge F_b(0,s,0)
       \ge F_b(0,2/b,0)
       =\frac{\beta+4q/b}
          {\beta(\beta+4q/b-4q/b^2)}.
\]
Here $2/b\le1$, since $b\ge4$.  The inequality that the final
expression be at most $b$ is, at $\beta=1/4$,
$b^2+8b-60\ge0$.  It contradicts $b<B_{\rm opp}$.
At $b=B_{\rm opp}$ the choice
$r=0$, $s=\gamma=2/b$, $\delta=0$, with every $a$-failure
routed to label $1$, attains equality.  This proves
\eqref{eq:r80-other-orientation} and completes the lemma.
\end{proof}

\subsection{A quantitative consequence of boundary rigidity}

The next estimate concerns the original, unrestricted six-parameter
controller.  It does not assume that its $a$-failure updates have
already been improved.  Write $J_0^{\rm rob}$ for its robust cost in
$\mathcal U_0$.

\begin{lemma}[Coercivity near the boundary optimizer]
\label{lem:r80-boundary-coercivity}
There are a neighborhood of the optimal controllers in
Lemma~\ref{lem:r80-boundary-rigidity} and constants $K,\Delta_0>0$
with the following property.  If a controller in that neighborhood
has $J_0^{\rm rob}\le B_0+\Delta$, where
$0\le\Delta\le\Delta_0$, then,
with $z=s+\delta$,
\begin{equation}\label{eq:r80-coercivity}
 z\le K\Delta,\qquad
 \gamma\le K\Delta z,\qquad
 |r-r_0|\le K\Delta.
\end{equation}
At $z=0$ the middle conclusion says $\gamma=0$.
\end{lemma}
\begin{proof}
Choose a small fixed neighborhood in which $r$ is bounded away
from zero and one and $s$ is smaller than $r/2$.
The actual model-$1$ cost is at least its cost after routing all
$a$-failures to label $2$.  In
\eqref{eq:r80-positive-polynomial}, the coefficients of
$\gamma,\delta^2,s\delta,s^2$ are bounded below by positive
constants on this neighborhood.  Also
$N_0$ is bounded above and below by positive constant multiples
of $z+\gamma$.  The denominator which converts that polynomial
to $C_1-f(D_0/N_0)$ is bounded above by a constant times $N_0$.
Thus, whenever $N_0>0$,
\begin{equation}\label{eq:r80-coercivity-core}
 C_1-f(1/C_0)\ge
 c\,\frac{\gamma+z^2}{z+\gamma}
\end{equation}
for some $c>0$.  Here $C_0$ is the actual baseline model-$0$
cost and $C_1$ may be the actual, unimproved model-$1$ cost.
Both are at most $B_0+\Delta$.  Since $f$ is increasing,
\[
 0\le C_1-f(1/C_0)
   \le B_0+\Delta-f\bigl(1/(B_0+\Delta)\bigr)
   \le K_1\Delta.
\]
The last bound follows from differentiability and
$f(1/B_0)=B_0$.

Shrink the neighborhood so that $z\le1$.  Then
$(\gamma+z^2)/(z+\gamma)\ge z$, and it is also at least
$\gamma/(z+\gamma)$.  These two bounds give
$z\le K_2\Delta$ and
$\gamma/(z+\gamma)\le K_2\Delta$.
For $K_2\Delta\le1/2$, the latter yields
$\gamma\le2K_2\Delta z$.

Put $t=1/C_0$.  The same inequalities give
$1/(B_0+\Delta)\le t\le f^{-1}(B_0+\Delta)$.
The derivative $f'(r_0)>0$, so $|t-r_0|\le K_3\Delta$.
Finally the explicit model-$0$ formula gives
\[
 r-t=\frac{(r-s)(1-r)\gamma}{N_0},
\]
which is bounded by a constant times
$\gamma/(z+\gamma)$.  This proves the estimate for $r$.
If $N_0=0$, then $z=\gamma=0$ and $t=r$ directly, so the
same conclusions hold.  Increasing one constant $K$ proves
all of \eqref{eq:r80-coercivity}.
\end{proof}

\subsection{The matching asymptotic}

\begin{proof}[Proof of Theorem~\ref{thm:r80-sharp-modulus}]
Every model has a successful command for $\varepsilon<1$.
Theorem~\ref{thm:r79-dark-feasibility} gives feasibility, and
Theorem~\ref{thm:r79-quantum-frontier} gives attainment at two
labels.  Alternatively, here the parameter cube
\eqref{eq:r80-six-parameters} is compact, total expected cost is
the supremum of its continuous finite-prefix costs, and a
constant mixed row gives a finite competitor.  These observations
prove attainment directly.

A one-label row with $a$-probability $r$ has costs
$1/[(1-\varepsilon)r]$ and $1/[\beta(1-r)]$.
Equating them proves \eqref{eq:r80-one-label-smooth}; this value
is at most $6$ when $\varepsilon\le1/2$.  The single fixed row
$r=1/3$ has cost at most $6$ for every model and every
$\varepsilon\in[0,1/2]$, so the fallback cap can be realized by
the same program throughout.  Pairing the unchanged model $1$
with itself and $0_0$ with $0_\varepsilon$, the scalar
nonterminal-map discrepancy is $\varepsilon$.  The cost
operators coincide.  The paired instrument discrepancy used in
Theorem~\ref{thm:r79-frontier-stability} is therefore exactly
$\varepsilon$.

For the upper estimate define the positive constants
\begin{equation}\label{eq:r80-leading-coefficients}
 A=r_0^{-2},\qquad
 F=\frac{\beta}{D^2}=f'(r_0),\qquad
 L=\frac{qr_0^2}{\beta D^2},\qquad
 k=\sqrt{\frac{F}{AL}}=\frac{\beta}{\sqrt q},\qquad
 h=\frac{k^{-1}-Lk}{A+F}.
\end{equation}
For sufficiently small $\varepsilon>0$, use the controller
\begin{equation}\label{eq:r80-asymptotic-controller}
 r=r_0+h\sqrt\varepsilon,\quad s=0,\quad
 \alpha=1,\quad\gamma=0,\quad
 \delta=k\sqrt\varepsilon,
\end{equation}
with any $\xi$.  These are legitimate probabilities.  Both
models are proper from every reachable label.  Direct substitution
in \eqref{eq:r80-two-state-cost} gives
\begin{align}
 C_0(\varepsilon)
   &=\frac1{(1-\varepsilon)r}
                 +\frac{\varepsilon}{(1-\varepsilon)\delta},
                       \label{eq:r80-upper-cost-zero}\\
 C_1(\varepsilon)
   &=\frac{\beta+q\delta+r}
        {\beta[\beta+qr+q(1-r)\delta]}.
                       \label{eq:r80-upper-cost-one}
\end{align}
Their expansions are respectively
\[
 B_0+(-Ah+k^{-1})\sqrt\varepsilon+O(\varepsilon),
 \qquad
 B_0+(Fh+Lk)\sqrt\varepsilon+O(\varepsilon).
\]
The two leading coefficients are equal by the definition of $h$,
and their common value is
\[
 \frac{2\sqrt{FAL}}{A+F}
        =\frac{2\sqrt q\,r_0^2}{D^2+\beta r_0^2}=c_0.
\]
This proves the required upper limit, in particular
$V_2(\varepsilon)-B_0=O(\sqrt\varepsilon)$.

For the converse take an attaining controller for each
$\varepsilon>0$.  For any fixed controller the model-$0$
cost is nondecreasing in $\varepsilon$.  Indeed the
nonterminal matrix is its matrix at zero plus
$\varepsilon$ times a nonnegative matrix, so each finite-prefix
cost is nondecreasing.  Model $1$ is unchanged.  Consequently
$V_2(\varepsilon)\ge B_0$, and every such optimizer has baseline
robust cost at most $V_2(\varepsilon)$.

Every limit point of these controllers as
$\varepsilon\downarrow0$ is a baseline optimizer.  To see this,
use compactness of the parameter cube and lower semicontinuity
of the baseline total cost, the latter being the supremum of
continuous finite-prefix costs.  The upper estimate just proved
then bounds a limiting baseline cost by $B_0$.
Lemma~\ref{lem:r80-boundary-rigidity} therefore implies
\[
 r\longrightarrow r_0,\quad s,\gamma,\delta\longrightarrow0,
                 \quad\alpha\longrightarrow1.
\]
No convergence of the immaterial coordinate $\xi$ is required.
For all sufficiently small parameters the coercivity lemma applies,
with $\Delta=V_2(\varepsilon)-B_0=O(\sqrt\varepsilon)$.
It gives
\begin{equation}\label{eq:r80-optimal-scales}
 s+\delta=O(\sqrt\varepsilon),\qquad
 \gamma=O(\varepsilon),\qquad
 r-r_0=O(\sqrt\varepsilon).
\end{equation}

Take any sequence along which the lower limit of the normalized
cost gap is realized.  Pass to a subsequence so that the bounded
scaled coordinates converge:
\[
 \frac{r-r_0}{\sqrt\varepsilon}\longrightarrow H,
 \qquad\frac{s}{\sqrt\varepsilon}\longrightarrow u,
 \qquad\frac{\delta}{\sqrt\varepsilon}\longrightarrow v,
 \qquad u,v\ge0.
\]
Write $t=u+v$.  We first show $t>0$.
In model $0_\varepsilon$ the two success probabilities and
off-diagonal nonterminal probabilities are
\begin{equation}\label{eq:r80-perturbed-chain}
 h_1=(1-\varepsilon)r,\quad h_2=(1-\varepsilon)s,
 \qquad
 t_1=\varepsilon r\alpha+(1-r)\gamma,\quad
 t_2=\varepsilon s\xi+(1-s)\delta.
\end{equation}
Let $\mathcal D=h_1h_2+h_1t_2+h_2t_1$.
For small positive $\varepsilon$, finite cost and
$\alpha\longrightarrow1$ require $s+\delta>0$;
otherwise the positive probability of entering the second label
leads to an absorbing failure state.  Thus $\mathcal D>0$, and
\eqref{eq:r80-two-state-cost} gives the exact identity
\begin{equation}\label{eq:r80-trap-cost}
 C_0(\varepsilon)-\frac1{h_1}
       =\frac{(h_1-h_2)t_1}{h_1\mathcal D}.
\end{equation}
In the neighborhood under consideration, $\mathcal D$ is at most
a constant times $s+\delta$, while $t_1\ge\varepsilon r\alpha$.
Hence the right-hand side is at least
$c\varepsilon/(s+\delta)$ for a fixed $c>0$.
Since $1/h_1=B_0+O(\sqrt\varepsilon)$ by
\eqref{eq:r80-optimal-scales}, the bounded normalized cost gap
excludes $(s+\delta)/\sqrt\varepsilon\longrightarrow0$.
Therefore $t>0$.

Now \eqref{eq:r80-perturbed-chain} and
\eqref{eq:r80-optimal-scales} give
$\mathcal D/\sqrt\varepsilon\longrightarrow r_0t$.
The nonnegative contribution of $\gamma$ to $t_1$ may be
dropped in a lower bound, while $\alpha\longrightarrow1$.
After dividing \eqref{eq:r80-trap-cost} by
$\sqrt\varepsilon$ we obtain
\begin{equation}\label{eq:r80-lower-zero}
 \liminf\frac{C_0(\varepsilon)-B_0}{\sqrt\varepsilon}
                       \ge-AH+\frac1t.
\end{equation}

For model $1$, route every $a$-failure to label $2$ only for the
purpose of a lower bound.  The resulting cost $N_1/D_1$ is a
smooth function of $r,s,\gamma,\delta$ in a neighborhood of the
baseline optimizer, since $D_1$ is strictly positive there.
Its derivatives with respect to $r,s,\delta$ at that point are
\[
 F,\qquad S=\frac{r_0}{\beta D},\qquad L,
 \qquad S-L=\frac{r_0}{D^2}>0.
\]
The $\gamma$ term is $O(\varepsilon)$ by
\eqref{eq:r80-optimal-scales}.  Taylor expansion, valid uniformly
on this fixed neighborhood, therefore yields
\begin{equation}\label{eq:r80-lower-one}
 \liminf\frac{C_1(\varepsilon)-B_0}{\sqrt\varepsilon}
       \ge FH+Su+Lv\ge FH+Lt.
\end{equation}
This bound applies to the original arbitrary $\alpha,\xi$;
the routing improvement only reduced its model-$1$ cost.

Combining \eqref{eq:r80-lower-zero} and
\eqref{eq:r80-lower-one}, the lower limit of the robust gap is
at least
\[
 \max\{-AH+t^{-1},\ FH+Lt\}
 \ge\frac{F/t+ALt}{A+F}
 \ge\frac{2\sqrt{FAL}}{A+F}=c_0.
\]
The first inequality is the convex combination with weights
$F/(A+F)$ and $A/(A+F)$; the second is the arithmetic--geometric
mean inequality.  Since the sequence realizing the lower limit
was arbitrary, this proves the lower limit in
\eqref{eq:r80-sharp-asymptotic}.  The upper construction gives the
matching upper limit.  Finally
$c_0\sqrt\varepsilon/\varepsilon^\alpha\longrightarrow\infty$
for every $\alpha>1/2$, proving the asserted optimality of the
exponent in the general primitive class.
\end{proof}

The limiting controller at zero perturbation is proper from its
prescribed start in both models, but its second label is dark and
unreachable in model $0_0$.  At a positive perturbation an
$a$-failure reaches that label with probability of order
$\varepsilon$.  An escape of order $\sqrt\varepsilon$ balances
the resulting residence cost against the cost of disturbing the
other model's optimal continuation.  The proof above quantifies
this explanation over all stochastic two-label controllers; it
does not restrict the converse to the displayed realizing program.
The exact leading constant belongs to this family.  The exponent
conclusion concerns the full class of
Theorem~\ref{thm:r79-frontier-stability}; it does not contradict
the Lipschitz result under the additional success floor in
Corollary~\ref{cor:r79-positive-success}.
